% =============================================================================
% Worked solutions for Simulados/mock2-midterm-topics-1-2b.pdf (Topics 1, 2A, 2B).
% Statements: mock2-midterm-topics-1-2b-statements.tex (shared with the exam paper).
% Compile: python -I Simulados/mock2-midterm-topics-1-2b-solutions_check.py (must print PASS),
%          then pdflatex (twice); check: pdffonts (Type 1 only, uni = yes).
% =============================================================================
\documentclass[11pt,a4paper]{article}
\usepackage[utf8]{inputenc}
\usepackage[T1]{fontenc}
\usepackage{lmodern}   % Type 1 vector fonts (replaces PK bitmaps)
\usepackage{cmap}      % ToUnicode CMap -> copyable text and accents
\usepackage[english]{babel}
\usepackage{amsmath,amssymb,amsfonts}
\usepackage[a4paper,margin=2.2cm]{geometry}
\usepackage{booktabs}
\usepackage{tabularx}
\usepackage{enumitem}
\usepackage{fancyhdr}
\usepackage{xcolor}
\usepackage{microtype}

\pagestyle{fancy}\fancyhf{}
\lhead{\footnotesize Finan\c{c}as II -- Corporativas (MPE Insper)}
\rhead{\footnotesize Mock midterm 2, Topics 1--2B --- Worked solutions}
\cfoot{\footnotesize\thepage}
\renewcommand{\headrulewidth}{0.3pt}
\setlength{\headheight}{14pt}

\makeatletter
\renewcommand{\@seccntformat}[1]{\csname the#1\endcsname.\quad}
\makeatother
\numberwithin{equation}{section}
\setlist{topsep=2pt,itemsep=1pt}
\newcolumntype{Y}{>{\raggedright\arraybackslash}X}

\newcommand{\srcline}[1]{\par\noindent{\small\color{black!65}\textbf{Source:} #1}\par\smallskip}
\newcommand{\statement}[1]{\par\smallskip\noindent\colorbox{black!5}{\parbox{\dimexpr\linewidth-2\fboxsep}{\small #1}}\par\medskip}
\newcommand{\step}[2]{\par\medskip\noindent\textbf{Step #1 --- #2.}~}
\newenvironment{notation}{\par\noindent\textbf{Notation.}\begin{itemize}[leftmargin=1.4em]}{\end{itemize}}
\newcommand{\intuition}[1]{\par\medskip\noindent\textbf{Intuition.}~#1}
\newcommand{\refline}[1]{\par\smallskip\noindent{\small\color{black!70}\textbf{Ref:} #1}}
\newcommand{\rubric}[2]{\par\smallskip\noindent{\small\textbf{Marking guide (#1 marks).}~#2}}
\newcommand{\answer}[1]{\par\medskip\noindent\fbox{\parbox{\dimexpr\linewidth-2\fboxsep-2\fboxrule}{\textbf{Answer.}~#1}}\par}
\newcommand{\TF}[3]{\item[#1] \textbf{#2.} #3}
\newcommand{\E}{\mathbb{E}}

% =============================================================================
% Statements of mock midterm 2 (Topics 1, 2A, 2B), written once and \input by
%   mock2-midterm-topics-1-2b.tex            (exam paper)
%   mock2-midterm-topics-1-2b-solutions.tex  (worked solutions)
% Every statement is a real exercise copied verbatim (rules/exercise-lists.md):
% source line first, then the full text. Allowed edits only: math in LaTeX, tables
% rebuilt, line breaks undone, currency written as \$.
% Needs: \srcline{...} defined by the including document.
% =============================================================================

% ---------------------------------------------------------------- Question 1.1
\newcommand{\QoneA}{%
\srcline{Higgins (Koski--Higgins--Mitton), \emph{Analysis for Financial Management} 13e, Ch.~1,
Exercise 3, p.~32 (PDF 53) $\cdot$ Answer in book: yes, pp.~397--398 (PDF 418--419)}

\noindent Answer true or false to each of the following. Briefly explain your reasoning for each answer.
\begin{enumerate}[label=\alph*.,leftmargin=2em]
  \item If a company gets into financial difficulty, it can use some of its shareholders' equity to pay
        its bills for a time.
  \item It is impossible for a firm to have a negative book value of equity without the firm going into
        bankruptcy.
  \item If a company increases its dividend, its net income will decrease.
  \item You can construct a sources and uses statement for 2020 if you have a company's balance sheets
        for year-end 2020 and 2021.
  \item The goodwill account on the balance sheet is an attempt by accountants to measure the benefits
        that result from a company's public relations efforts in the community.
  \item A reduction in an asset account is a source of cash, while a reduction in a liability account is
        a use of cash.
  \item A company's market value of equity must always be higher than its book value of equity.
\end{enumerate}}

% ---------------------------------------------------------------- Question 1.2
\newcommand{\QoneB}{%
\srcline{Higgins (Koski--Higgins--Mitton), \emph{Analysis for Financial Management} 13e, Ch.~8,
Exercise 1, p.~327 (PDF 348) $\cdot$ Answer in book: yes, pp.~421--422 (PDF 442--443)}

\noindent Is each of the following statements true or false? Explain your answers briefly.
\begin{enumerate}[label=\alph*.,leftmargin=2em]
  \item Using the same risk-adjusted discount rate to discount all future cash flows ignores the fact
        that the more distant cash flows are often riskier than cash flows occurring sooner.
  \item The cost of capital, or WACC, is not the correct discount rate to use for all projects
        undertaken by a firm.
  \item If you can borrow all of the money you need for a project at 6 percent, the cost of capital for
        this project is 6 percent.
  \item The best way to estimate the cost of debt capital for a firm is to divide the interest expense
        on the income statement by the interest-bearing debt on the balance sheet.
  \item One reliable estimate of a privately held firm's equity beta is the average of the equity betas
        of several publicly held competitors.
  \item Nominal cash flows should always be discounted using a nominal discount rate.
\end{enumerate}}

% ---------------------------------------------------------------- Question 2.1
\newcommand{\QtwoA}{%
\srcline{Berk, DeMarzo \& Harford, \emph{Fundamentals of Corporate Finance} 2e, Ch.~13, Exercises 13.1
and 13.3, p.~403 (PDF 434) $\cdot$ Answer in book: no}

\noindent\textbf{1.} MV Corporation has debt with market value of \$100 million, common equity with a
book value of \$100 million, and preferred stock worth \$20 million outstanding. Its common equity trades
at \$50 per share, and the firm has 6 million shares outstanding. What weights should MV Corporation use
in its WACC?

\smallskip\noindent\textbf{3.} Book Co. has 1 million shares of common equity with a par (book) value of
\$1, retained earnings of \$30 million, and its shares have a market value of \$50 per share. It also has
debt with a par value of \$20 million that is trading at 101\% of par.
\begin{enumerate}[label=\alph*.,leftmargin=2em]
  \item What is the market value of its equity?
  \item What is the market value of its debt?
  \item What weights should it use in computing its WACC?
\end{enumerate}}

% ---------------------------------------------------------------- Question 2.2
\newcommand{\QtwoB}{%
\srcline{Berk, DeMarzo \& Harford, \emph{Fundamentals of Corporate Finance} 2e, Ch.~13, Exercises
13.18--13.19, p.~405 (PDF 436) $\cdot$ Answer in book: no}

\noindent\textbf{18.} RiverRocks (whose WACC is 12\%) is considering an acquisition of Raft Adventures
(whose WACC is 15\%). What is the appropriate discount rate for RiverRocks to use to evaluate the
acquisition? Why?

\smallskip\noindent\textbf{19.} RiverRocks' purchase of Raft Adventures (from Problem 18) will cost \$100
million, but will generate cash flows that start at \$15 million in one year and then grow at 4\% per year
forever. What is the NPV of the acquisition?}

% ---------------------------------------------------------------- Question 2.3
\newcommand{\QtwoC}{%
\srcline{Berk, DeMarzo \& Harford, \emph{Fundamentals of Corporate Finance} 2e, Ch.~13, Exercise 13.15,
p.~404 (PDF 435) $\cdot$ Answer in book: no}

\noindent Growth Company's current share price is \$20 and it is expected to pay a \$1 dividend per share
next year. After that, the firm's dividends are expected to grow at a rate of 4\% per year.
\begin{enumerate}[label=\alph*.,leftmargin=2em]
  \item What is an estimate of Growth Company's cost of equity?
  \item Growth Company also has preferred stock outstanding that pays a \$2 per share fixed dividend. If
        this stock is currently priced at \$28, what is Growth Company's cost of preferred stock?
  \item Growth Company has existing debt issued three years ago with a coupon rate of 6\%. The firm just
        issued new debt at par with a coupon rate of 6.5\%. What is Growth Company's pre-tax cost of
        debt?
  \item Growth Company has 5 million common shares outstanding and 1 million preferred shares
        outstanding, and its equity has a total book value of \$50 million. Its liabilities have a
        market value of \$20 million. If Growth Company's common and preferred shares are priced as in
        parts (a) and (b), what is the market value of Growth Company's assets?
  \item Growth Company faces a 35\% tax rate. Given the information in parts (a)--(d), and your answers
        to those problems, what is Growth Company's WACC?
\end{enumerate}}

% ---------------------------------------------------------------- Question 2.4
\newcommand{\QtwoD}{%
\srcline{Berk, DeMarzo \& Harford, \emph{Fundamentals of Corporate Finance} 2e, Ch.~16, Exercise 16.3,
p.~491 (PDF 522); section instruction on p.~490 (PDF 521) $\cdot$ Answer in book: no}

\noindent\emph{Section instruction (Capital Structure in Perfect Capital Markets):} For problems in this
section assume no taxes or distress costs.

\smallskip\noindent Acort Industries owns assets that will have an 80\% probability of having a market
value of \$50 million in one year. There is a 20\% chance that the assets will be worth only \$20
million. The current risk-free rate is 5\%, and Acort's assets have a cost of capital of 10\%.
\begin{enumerate}[label=\alph*.,leftmargin=2em]
  \item If Acort is unlevered, what is the current market value of its equity?
  \item Suppose instead that Acort has debt with a face value of \$20 million due in one year. According
        to MM, what is the value of Acort's equity in this case?
  \item What is the expected return of Acort's equity without leverage? What is the expected return of
        Acort's equity with leverage?
  \item What is the lowest possible realized return of Acort's equity with and without leverage?
\end{enumerate}}

% ---------------------------------------------------------------- Question 3.1
\newcommand{\QthreeA}{%
\srcline{Berk, DeMarzo \& Harford, \emph{Fundamentals of Corporate Finance} 2e, Ch.~16, Exercise 16.10,
p.~492 (PDF 523) $\cdot$ Answer in book: no}

\noindent Suppose the corporate tax rate is 40\%. Consider a firm that earns \$1000 before interest and
taxes each year with no risk. The firm's capital expenditures equal its depreciation expenses each year,
and it will have no changes to its net working capital. The risk-free interest rate is 5\%.
\begin{enumerate}[label=\alph*.,leftmargin=2em]
  \item Suppose the firm has no debt and pays out its net income as a dividend each year. What is the
        value of the firm's equity?
  \item Suppose instead the firm makes interest payments of \$500 per year. What is the value of equity?
        What is the value of debt?
  \item What is the difference between the total value of the firm with leverage and without leverage?
  \item To what percentage of the value of the debt is the difference in part (c) equal?
\end{enumerate}}

% ---------------------------------------------------------------- Question 3.2
\newcommand{\QthreeB}{%
\srcline{Berk, DeMarzo \& Harford, \emph{Fundamentals of Corporate Finance} 2e, Ch.~16, Exercises
16.16--16.17, p.~493 (PDF 524) $\cdot$ Answer in book: no}

\noindent\textbf{16.} Milton Industries expects free cash flows of \$5 million each year. Milton's
corporate tax rate is 35\%, and its unlevered cost of capital is 15\%. The firm also has outstanding debt
of \$19.05 million, and it expects to maintain this level of debt permanently.
\begin{enumerate}[label=\alph*.,leftmargin=2em]
  \item What is the value of Milton Industries without leverage?
  \item What is the value of Milton Industries with leverage?
\end{enumerate}
\noindent\textbf{17.} NatNah, a builder of acoustic accessories, has no debt and an equity cost of
capital of 15\%. NatNah decides to increase its leverage to maintain a market debt-to-value ratio of 0.5.
Suppose its debt cost of capital is 9\% and its corporate tax rate is 35\%. If NatNah's pretax WACC
remains constant, what will its (effective after-tax) WACC be with the increase in leverage?}

% ---------------------------------------------------------------- Question 3.3
\newcommand{\QthreeC}{%
\srcline{Berk, DeMarzo \& Harford, \emph{Fundamentals of Corporate Finance} 2e, Ch.~16, Exercise 16.23,
p.~494 (PDF 525) $\cdot$ Answer in book: no}

\noindent Marpor Industries has no debt and expects to generate free cash flows of \$16 million each year.
Marpor believes that if it permanently increases its level of debt to \$40 million, the risk of financial
distress may cause it to lose some customers and receive less favorable terms from its suppliers. As a
result, Marpor's expected free cash flows with debt will be only \$15 million per year. Suppose Marpor's
tax rate is 35\%, the risk-free rate is 5\%, the expected return of the market is 15\%, and the beta of
Marpor's free cash flows is 1.10 (with or without leverage).
\begin{enumerate}[label=\alph*.,leftmargin=2em]
  \item Estimate Marpor's value without leverage.
  \item Estimate Marpor's value with the new leverage.
\end{enumerate}}

% ---------------------------------------------------------------- Question 3.4
\newcommand{\QthreeD}{%
\srcline{Berk, DeMarzo \& Harford, \emph{Fundamentals of Corporate Finance} 2e, Ch.~17, Exercise 17.9,
p.~528 (PDF 559) $\cdot$ Answer in book: no}

\noindent The HNH Corporation will pay a constant dividend of \$2 per share, per year, in perpetuity.
Assume all investors pay a 20\% tax on dividends and that there is no capital gains tax. The cost of
capital for investing in HNH stock is 12\%.
\begin{enumerate}[label=\alph*.,leftmargin=2em]
  \item What is the price of a share of HNH stock?
  \item Assume that management makes a surprise announcement that HNH will no longer pay dividends but
        will use the cash to repurchase stock instead. What is the price of a share of HNH stock now?
\end{enumerate}}


\begin{document}

\begin{center}
  {\Large\bfseries Finance II --- Mock Midterm Exam, Set 2 (Topics 1, 2A, 2B)}\\[4pt]
  {\large Worked solutions}
\end{center}

{\sloppy\noindent Exam paper: \texttt{Simulados/mock2-midterm-topics-1-2b.pdf}. Each question's full
statement (grey box, with its source) comes first, then the solution, a boxed answer and a marking
guide. Money is in the units of each statement. Rates are decimals inside equations. Every number is
reproduced by \texttt{Simulados/mock2-midterm-topics-1-2b-solutions\_check.py}. For the true/false items
Higgins prints suggested answers; each verdict below agrees with the book's, and the reason is
written in this course's terms. References: lecture slides \texttt{Aula\_Topic1}, \texttt{Aula\_Topic2A},
\texttt{Aula\_Topic2B}; formula sheet \texttt{Notas/formula-sheet-topics-1-2b.pdf}.\par}

% =============================================================================
\section{Question 1.1 --- The balance sheet and cash (true/false)}

\statement{\QoneA}

\begin{enumerate}[leftmargin=2.2em,itemsep=4pt]
  \TF{a.}{False}{Equity is a \emph{claim} on the right-hand side of the balance sheet, not a pool of money.
  The cash shareholders contributed, plus retained profits, has already been spent on assets. Bills are
  paid with cash, which is an asset.}
  \TF{b.}{False}{Book equity is the plug $A-L$. Accumulated losses or large buybacks can make it
  negative while the firm keeps paying its bills and its market equity is large (Starbucks 2020:
  $-7.9$ bn book vs.\ $122.6$ bn market). Bankruptcy is about failing to pay obligations when due.}
  \TF{c.}{False}{Dividends are a distribution of net income after it is computed. A higher dividend
  lowers retained earnings, not net income.}
  \TF{d.}{False}{A sources and uses statement for 2020 compares the balance sheet at the start of 2020
  (year-end 2019) with the one at the end of 2020. Year-end 2020 and 2021 give the statement for 2021.}
  \TF{e.}{False}{Goodwill records the amount an acquirer paid above the book value of the acquired
  firm's assets. It has nothing to do with public relations.}
  \TF{f.}{True}{Shrinking an asset (collecting receivables, selling inventory or equipment) frees cash;
  shrinking a liability (paying suppliers, repaying debt) consumes cash. These are the same signs as in
  the indirect cash flow statement: $-\Delta AR$, $+\Delta AP$.}
  \TF{g.}{False}{Market equity usually exceeds book equity (intangibles and growth are not on the
  balance sheet), but a firm with obsolete assets or a failing strategy can trade below book
  (market-to-book $<1$).}
\end{enumerate}

\answer{a F, b F, c F, d F, e F, f T, g F (the book's answers, pp.~397--398).}
\refline{Topic 1 slides (book vs.\ market balance sheet, accruals, cash flow statement).}
\rubric{14}{2 per item: 1 for the verdict, 1 for a reason that names the mechanism (claim vs.\ asset,
plug, retained earnings, timing of balance sheets, acquisition premium, signs, market-to-book).}

% =============================================================================
\section{Question 1.2 --- The cost of capital (true/false)}

\statement{\QoneB}

\begin{enumerate}[leftmargin=2.2em,itemsep=4pt]
  \TF{a.}{False}{With a constant rate, $PV=CF_t/(1+r)^t$: the risk premium inside $r$ compounds with
  $t$, so later cash flows are already discounted more for risk. A single rate assumes risk grows at a
  constant geometric rate.}
  \TF{b.}{True}{The WACC is right only for projects with the risk (and financing) of the firm's existing
  assets. A riskier or safer project needs its own rate, taken from comparables (Question~2.2 is an
  example).}
  \TF{c.}{False}{This is the marginal-cost-of-capital fallacy. Borrowing for the project raises the risk
  borne by equity, so $r_E$ rises (MM~II). The discount rate is the opportunity cost of the project's
  risk, not the cost of whichever source pays for it.}
  \TF{d.}{False}{Interest expense over book debt is a historical coupon rate, on a debt balance that
  changes during the year. The cost of debt is the return lenders require \emph{now}: the yield to
  maturity on the firm's debt (and, with default risk, the expected return is below even that).}
  \TF{e.}{False}{Competitors' equity betas mix business risk with their own leverage. Unlever each
  ($\beta_A=\frac EV\beta_E+\frac DV\beta_D$), average the asset betas, then relever at the private
  firm's own debt ratio.}
  \TF{f.}{True}{Discount like with like: nominal cash flows at nominal rates, real cash flows at real
  rates. Mixing them double counts or ignores inflation.}
\end{enumerate}

\answer{a F, b T, c F, d F, e F, f T (the book's answers, pp.~421--422).}
\refline{Topic 2A slides (project discount rate, unlevering comparables); Higgins Ch.~8.}
\rubric{12}{2 per item: 1 for the verdict, 1 for the reason; (c) and (e) need the leverage argument.}

% =============================================================================
\section{Question 2.1 --- Market-value weights}

\statement{\QtwoA}

\step{1}{MV Corporation} The WACC weights are market values. The book value of common equity
(\$100 million) is irrelevant. Market value of equity: $6\times50=300$; total $V=100+300+20=420$:
\[
  w_E=\frac{300}{420}=0.7143,\qquad w_D=\frac{100}{420}=0.2381,\qquad w_P=\frac{20}{420}=0.0476 .
\]
\step{2}{Book Co.} (a) $E=1\times50=\$50$ million (book equity would be $1+30=\$31$ million: not used).
(b) $D=20\times1.01=\$20.2$ million. (c) $V=70.2$; $w_E=50/70.2=0.7123$, $w_D=20.2/70.2=0.2877$.

\answer{MV: 71.43\% common equity, 23.81\% debt, 4.76\% preferred. Book Co.: (a) \$50 million; (b) \$20.2
million; (c) 71.23\% equity, 28.77\% debt.}
\intuition{The weights describe how much investors have at stake today, which is what they require a
return on. Book values describe what was paid in the past.}
\refline{Topic 1 slides (book vs.\ market value); Topic 2A (WACC).}
\rubric{6}{3 for MV (market equity computed, book value ignored); 3 for Book Co.\ (1 each part).}

% =============================================================================
\section{Question 2.2 --- The project's own discount rate}

\statement{\QtwoB}

\step{1}{Which rate (Problem 18)} 15\%, Raft Adventures' WACC. The discount rate belongs to the
cash flows being valued: the acquired business keeps its own risk whoever owns it. RiverRocks' 12\%
reflects RiverRocks' assets, not Raft's.

\step{2}{NPV (Problem 19)} A growing perpetuity starting in one year:
\[
  PV=\frac{C_1}{r-g}=\frac{15}{0.15-0.04}=136.36,\qquad NPV=136.36-100=\$36.36\text{ million}.
\]
\textbf{Contrast.} At RiverRocks' 12\% the NPV would be $15/0.08-100=\$87.5$ million: the wrong rate
more than doubles the NPV and would justify overpaying.

\answer{18: 15\%, the target's cost of capital, because it reflects the risk of the cash flows acquired.
19: $NPV=\$36.36$ million.}
\refline{Topic 2A slides (project risk, growing perpetuity).}
\rubric{8}{3 for 15\% with the reason; 4 for the NPV ($C_1/(r-g)$ with $C_1=15$); 1 for a sensible
comment.}

% =============================================================================
\section{Question 2.3 --- Estimating a WACC}

\statement{\QtwoC}

\step{1}{Cost of equity (a)} Constant-growth dividend model: $r_E=\frac{Div_1}{P_0}+g=\frac{1}{20}+0.04=9\%$.

\step{2}{Cost of preferred (b)} A perpetuity: $r_P=\frac{2}{28}=7.14\%$.

\step{3}{Cost of debt (c)} The relevant cost is what new debt costs today: the new issue at par
pays 6.5\%, so its yield is 6.5\%. The 6\% coupon from three years ago is history.

\step{4}{Market value of assets (d)} Market value of the claims: common $5\times20=100$, preferred
$1\times28=28$, liabilities $20$; total $\$148$ million. The book value of equity (\$50 million) is not
used.

\step{5}{WACC (e)}
\[
  r_{WACC}=\tfrac{100}{148}(9\%)+\tfrac{28}{148}(7.14\%)+\tfrac{20}{148}(6.5\%)(1-0.35)
  =6.08\%+1.35\%+0.57\%=8.00\% .
\]
Only debt gets the $(1-\tau)$: interest is deductible, preferred dividends are not.

\answer{(a) 9\%; (b) 7.14\%; (c) 6.5\%; (d) \$148 million; (e) $r_{WACC}=8.00\%$.}
\refline{Topic 2A slides (WACC); Berk Ch.~13.}
\rubric{10}{2 each part; (e) needs market weights and the after-tax cost of debt only.}

% =============================================================================
\section{Question 2.4 --- MM and the risk of levered equity}

\statement{\QtwoD}

\step{1}{Unlevered equity (a)} $\E[V_1]=0.8(50)+0.2(20)=44$; $V_0=44/1.10=\$40$ million.

\step{2}{Levered equity (b)} The debt pays \$20 million in both states (assets are never below 20),
so it is risk-free: $D=20/1.05=19.05$. By MM~I the firm is still worth 40, so $E=40-19.05=\$20.95$
million.

\step{3}{Expected returns (c)} Unlevered: 10\% (the assets' cost of capital). Levered: equity pays
$50-20=30$ or $20-20=0$, so $\E[\text{payoff}]=0.8(30)=24$ and
\[
  r_E=\frac{24}{20.95}-1=14.55\% .
\]
Check with MM~II: $r_E=r_A+\frac DE(r_A-r_D)=10\%+\frac{19.05}{20.95}(10\%-5\%)=14.55\%$.

\step{4}{Lowest realized returns (d)} Without leverage: $20/40-1=-50\%$. With leverage:
$0/20.95-1=-100\%$ (the best case is $+25\%$ unlevered and $+43.18\%$ levered).

\answer{(a) \$40 million; (b) \$20.95 million; (c) 10\% unlevered, 14.55\% levered; (d) $-50\%$
unlevered, $-100\%$ levered.}
\intuition{Leverage does not change the value of the firm, but it concentrates the same asset risk on a
smaller equity claim: the expected return rises because the downside is worse, not because value is
created.}
\refline{Topic 2A slides (MM~I and II, homemade leverage).}
\rubric{10}{2 for (a); 3 for (b) with safe debt at $r_f$; 3 for (c) with the MM~II check; 2 for (d).}

% =============================================================================
\section{Question 3.1 --- The value of the interest tax shield}

\statement{\QthreeA}

\step{1}{No debt (a)} Net income $=1000(1-0.4)=600$ a year, riskless, paid as dividends forever:
$E=600/0.05=\$12{,}000$.

\step{2}{With \$500 interest (b)} Net income $=(1000-500)(0.6)=300$, so $E=300/0.05=\$6{,}000$. Debt
pays a riskless \$500 forever: $D=500/0.05=\$10{,}000$.

\step{3}{Difference (c)} $V_L=6{,}000+10{,}000=16{,}000$; $V_L-V_U=16{,}000-12{,}000=\$4{,}000$.

\step{4}{As a share of debt (d)} $4{,}000/10{,}000=40\%=\tau$. The interest tax shield is
$0.4\times500=200$ a year, worth $200/0.05=4{,}000=\tau D$.

\answer{(a) \$12,000; (b) equity \$6,000, debt \$10,000; (c) \$4,000; (d) 40\% of the debt, the tax
rate: $V_L=V_U+\tau D$ for permanent debt.}
\intuition{Investors in total receive $1000-0.4(1000-500)=800$ a year with debt vs.\ 600 without: the
government's take falls by 200, and that is the whole gain.}
\refline{Topic 2B slides (interest tax shield, permanent debt).}
\rubric{10}{2 for (a); 4 for (b); 2 for (c); 2 for (d) linked to $\tau D$.}

% =============================================================================
\section{Question 3.2 --- Fixed debt vs.\ a constant debt-to-value ratio}

\statement{\QthreeB}

\step{1}{Milton (16): permanent debt} (a) $V_U=5/0.15=\$33.33$ million. (b) Permanent debt
$\Rightarrow$ shields as safe as the debt, $PV(ITS)=\tau D=0.35(19.05)=6.67$:
$V_L=33.33+6.67=\$40.00$ million.

\step{2}{NatNah (17): constant $D/V$} The pretax WACC stays at $r_U=15\%$. After tax:
\[
  r_{WACC}=r_U-\tau r_D\frac DV=15\%-0.35(9\%)(0.5)=15\%-1.575\%=13.425\% .
\]
Check: $15\%=0.5\,r_E+0.5(9\%)\Rightarrow r_E=21\%$; then $0.5(21\%)+0.5(9\%)(0.65)=13.425\%$.

\answer{16: (a) \$33.33 million; (b) \$40.00 million. 17: after-tax WACC $=13.425\%$ (13.43\%).}
\intuition{Two debt policies, two tools: a fixed amount of debt is valued by APV with $PV(ITS)=\tau D$;
a constant $D/V$ goes into the WACC, which falls below $r_U$ by exactly $\tau r_DD/V$.}
\refline{Topic 2B slides (debt-policy table); cheatsheet §3.}
\rubric{10}{4 for Milton ($V_U$, then $+\tau D$); 6 for NatNah (pretax WACC $=r_U$, then subtract
$\tau r_DD/V$ or rebuild from $r_E$).}

% =============================================================================
\section{Question 3.3 --- Tax shield vs.\ financial distress}

\statement{\QthreeC}

\step{1}{Discount rate} The beta of the free cash flows gives the unlevered cost of capital:
$r_U=5\%+1.10(15\%-5\%)=16\%$.

\step{2}{Without leverage (a)} $V_U=16/0.16=\$100$ million.

\step{3}{With leverage (b)} The expected FCF fall to 15 (indirect distress costs); the permanent debt
adds $PV(ITS)=\tau D=0.35(40)=14$:
\[
  V_L=\frac{15}{0.16}+0.35(40)=93.75+14=\$107.75\text{ million}.
\]
In APV form: $V_L=V_U+PV(ITS)-PV(\text{distress})=100+14-6.25$.

\answer{(a) \$100 million; (b) \$107.75 million: the tax shield (14) exceeds the PV of distress costs
(6.25), so the new leverage adds \$7.75 million.}
\intuition{This is the static trade-off with numbers: leverage pays while the marginal shield exceeds
the marginal expected distress cost.}
\refline{Topic 2B slides ($V_L=V_U+PV(ITS)-PV(\text{distress})$, indirect costs).}
\rubric{10}{3 for $r_U$ by CAPM; 2 for (a); 5 for (b): lower FCF discounted at $r_U$ (2), $+\tau D$
(2), conclusion (1).}

% =============================================================================
\section{Question 3.4 --- Payout policy and personal taxes}

\statement{\QthreeD}

\begin{notation}
  \item The 12\% is read as the \emph{after-tax} return investors require on HNH (what they earn on
        equally risky stock); capital gains are untaxed.
\end{notation}

\step{1}{With dividends (a)} Investors keep $2(1-0.2)=1.60$ a year:
$P=\frac{1.60}{0.12}=\$13.33$.

\step{2}{With repurchases (b)} The same \$2 a share now reaches investors as capital gains, which are
untaxed: $P=\frac{2}{0.12}=\$16.67$, a rise of \$3.33 (25\%) on the announcement.

\answer{(a) \$13.33; (b) \$16.67.}
\intuition{With a dividend tax and no capital gains tax, paying out by repurchase saves investors the
tax, and the saving is capitalised in the price at once. In a perfect market both prices would be
equal.}
\textbf{Watch.} Reading 12\% as a pre-tax rate gives \$16.67 in (a) and hides the point of the
question, which is the tax disadvantage of dividends.
\refline{Topic 2B slides (payout with taxes); Berk §17.3.}
\rubric{10}{5 for (a) with the after-tax dividend; 5 for (b) with the reason (capital gains untaxed).}

\end{document}
